% GENERATED FILE. Edit canonical lessons, then run npm run generate:books.
% canonical-source-hash: fnv1a64:2a0449c46a95b7f2
% canonical-lessons: TE-C29-subhodayam

\chapter{Good Morning}
\label{ch:te-good-morning}

This chapter is generated from the canonical micro-lessons used by Language Ladder.
Each section stays independently resumable and preserves the authored prerequisite order.

\section[\te{శుభోదయం}]{\te{శుభోదయం} (śubhōdayam) — a claim worth real skepticism, not a settled fact}

\label{lesson:TE-C29-subhodayam}



\begin{quote}
\textbf{Warm-up.} {[}PAUSE 2s{]} You already have \te{ఉదయం}. This lesson's honest twist is about a dramatic-sounding claim — and about being skeptical of it, even though it came from a source fetched directly.
\end{quote}

\begin{cousinweb}
\textbf{\te{శుభోదయం}} (\textbf{śubhōdayam}) — "\textbf{good morning}" — is built from \textbf{\te{శుభ}} ("good") plus \textbf{\te{ఉదయం}} (already met, Chapter 26). Grammatically and etymologically, this is the straightforward, traditional Telugu "good morning."
\end{cousinweb}

\begin{culture}
Here's the twist, and the honesty this time cuts against the source itself: one directly-fetched page describes \textbf{\te{శుభోదయం}} as \textbf{out of trend}, displaced by \textbf{English} "good morning," with children taught the English phrase from \textbf{preschool}. That's a vivid, quotable claim — but "fetched directly" doesn't make a source authoritative. This is the \textbf{same genre} of Telugu/Kannada/Hindi-learning content blog already found \textbf{overstated} for Hindi's \dv{सुप्रभात} and Kannada's own \textbf{\ka{ಶುಭೋದಯ}} (a single uncredentialed listicle, no sociolinguistic methodology, no citations) — treat this claim the same skeptical way, not as settled.

Genuinely worth separating from that dramatic claim: \textbf{\te{సుప్రభాతం}} (\emph{suprabhātam}) also exists — built on \textbf{\te{ప్రభాత}} (\emph{prabhāta}), a \textbf{completely different} Sanskrit root from \textbf{\te{ఉదయ}} (\emph{udaya}). This is \textbf{not} the same root behind Kannada's \textbf{\ka{ಶುಭೋದಯ}}, which itself uses \emph{udaya} — don't conflate the two Sanskrit "morning" roots. A third option, \textbf{\te{ప్రొద్దుటి} \te{వందనాలు}} ("morning greetings," built on \textbf{\te{పొద్దు}}, Chapter 26's genuine native competitor to \te{ఉదయం}), also exists but isn't confirmed common either. \textbf{\te{నమస్కారం}} is the safer bet for everyday use — matching the same \te{నమస్కార}/\dv{नमस्ते} pattern already found doing the real everyday work for both Kannada and Hindi.
\end{culture}

\subsection*{Guided Practice}
{[}PAUSE 1s{]}

\begin{itemize}
\raggedright
  \item {[}YOU SAY: "śubhōdayam" — "good morning," śubha + udayam{]}
  \item {[}YOU SAY: the honest twist — one low-authority source claims it's out of trend, displaced by English; treat with real skepticism{]}
  \item {[}YOU SAY: \te{నమస్కారం} — the safer everyday-greeting bet{]}
\end{itemize}

\begin{tcolorbox}[breakable,colback=teal!4,colframe=teal!35!black,title={Wrap-up recall}]
{[}PAUSE 3s{]} What is \te{శుభోదయం} built from? (\textbf{\te{శుభ}} + \textbf{\te{ఉదయం}}.) Is the claim that \te{శుభోదయం} has been displaced by English a confirmed, well-sourced fact? (\textbf{No} — it comes from one low-authority content-blog source, the same genre already found overstated for Hindi/Kannada; treat it skeptically, not as settled.) Does \te{సుప్రభాతం} share \te{ఉదయం}'s root, or Kannada's \ka{ಶುಭೋದಯ}'s root? (\textbf{Neither} — its own root, \te{ప్రభాత}/prabhāta, is a genuinely different Sanskrit word.)
\end{tcolorbox}
